\documentclass[10pt,a4paper]{amsart}
\usepackage{amsmath,amssymb,amsthm,graphicx}
\usepackage{verbatim}
\newenvironment{reference}{\comment}{\endcomment}
%\newenvironment{reference}{}{}
\newenvironment{slogan}{\comment}{\endcomment}
\newenvironment{history}{\comment}{\endcomment}

% For commutative diagrams we use Xy-pic
\usepackage[all]{xy}

% We use 2cell for 2-commutative diagrams.
\xyoption{2cell}
% Theorem environments.
%
\theoremstyle{plain}
\newtheorem{theorem}[subsection]{Theorem}
\newtheorem{proposition}[subsection]{Proposition}
\newtheorem{lemma}[subsection]{Lemma}

\theoremstyle{definition}
\newtheorem{definition}[subsection]{Definition}
\newtheorem{example}[subsection]{Example}
\newtheorem{exercise}[subsection]{Exercise}
\newtheorem{situation}[subsection]{Situation}

\theoremstyle{remark}
\newtheorem{remark}[subsection]{Remark}
\newtheorem{remarks}[subsection]{Remarks}

\numberwithin{equation}{subsection}

% Macros
%
\def\lim{\mathop{\mathrm{lim}}\nolimits}
\def\colim{\mathop{\mathrm{colim}}\nolimits}
\def\Spec{\mathop{\mathrm{Spec}}}
\def\Hom{\mathop{\mathrm{Hom}}\nolimits}
\def\Ext{\mathop{\mathrm{Ext}}\nolimits}
\def\SheafHom{\mathop{\mathcal{H}\!\mathit{om}}\nolimits}
\def\SheafExt{\mathop{\mathcal{E}\!\mathit{xt}}\nolimits}
\def\Sch{\mathit{Sch}}
\def\Mor{\mathop{\mathrm{Mor}}\nolimits}
\def\Ob{\mathop{\mathrm{Ob}}\nolimits}
\def\Sh{\mathop{\mathit{Sh}}\nolimits}
\def\NL{\mathop{N\!L}\nolimits}
\def\CH{\mathop{\mathrm{CH}}\nolimits}
\def\proetale{{pro\text{-}\acute{e}tale}}
\def\etale{{\acute{e}tale}}
\def\QCoh{\mathit{QCoh}}
\def\Ker{\mathop{\mathrm{Ker}}}
\def\Im{\mathop{\mathrm{Im}}}
\def\Coker{\mathop{\mathrm{Coker}}}
\def\Coim{\mathop{\mathrm{Coim}}}

% Boxtimes
%
\DeclareMathSymbol{\boxtimes}{\mathbin}{AMSa}{"02}

%
% Macros for moduli stacks/spaces
%
\def\QCohstack{\mathcal{QC}\!\mathit{oh}}
\def\Cohstack{\mathcal{C}\!\mathit{oh}}
\def\Spacesstack{\mathcal{S}\!\mathit{paces}}
\def\Quotfunctor{\mathrm{Quot}}
\def\Hilbfunctor{\mathrm{Hilb}}
\def\Curvesstack{\mathcal{C}\!\mathit{urves}}
\def\Polarizedstack{\mathcal{P}\!\mathit{olarized}}
\def\Complexesstack{\mathcal{C}\!\mathit{omplexes}}
% \Pic is the operator that assigns to X its picard group, usage \Pic(X)
% \Picardstack_{X/B} denotes the Picard stack of X over B
% \Picardfunctor_{X/B} denotes the Picard functor of X over B
\def\Pic{\mathop{\mathrm{Pic}}\nolimits}
\def\Picardstack{\mathcal{P}\!\mathit{ic}}
\def\Picardfunctor{\mathrm{Pic}}
\def\Deformationcategory{\mathcal{D}\!\mathit{ef}}
\usepackage{cleveref}
\crefname{originallemma}{original lemma}{original lemmas}
\Crefname{originallemma}{Original lemma}{Original lemmas}
\makeatletter
\newlabel{duality-lemma-equivalent-definitions}{{0A86}{source}{}{sourcepointer-1}{}}
\newlabel{duality-lemma-equivalent-definitions@cref}{{[originallemma][0][]0A86}{[1][0][]source}{}{}{}}
\makeatother
\crefname{originallemma}{original lemma}{original lemmas}
\Crefname{originallemma}{Original lemma}{Original lemmas}
\makeatletter
\newlabel{dualizing-lemma-dualizing}{{0A7C}{source}{}{sourcepointer-2}{}}
\newlabel{dualizing-lemma-dualizing@cref}{{[originallemma][0][]0A7C}{[1][0][]source}{}{}{}}
\makeatother
\crefname{originalreference}{original reference}{original references}
\Crefname{originalreference}{Original reference}{Original references}
\makeatletter
\newlabel{lemma-affine-duality}{{0A88}{source}{}{sourcepointer-3}{}}
\newlabel{lemma-affine-duality@cref}{{[originalreference][0][]0A88}{[1][0][]source}{}{}{}}
\makeatother
\crefname{originalreference}{original reference}{original references}
\Crefname{originalreference}{Original reference}{Original references}
\makeatletter
\newlabel{lemma-equivalent-definitions}{{0A86}{source}{}{sourcepointer-4}{}}
\newlabel{lemma-equivalent-definitions@cref}{{[originalreference][0][]0A86}{[1][0][]source}{}{}{}}
\makeatother
\crefname{originallemma}{original lemma}{original lemmas}
\Crefname{originallemma}{Original lemma}{Original lemmas}
\makeatletter
\newlabel{sites-cohomology-lemma-internal-hom-evaluate}{{0A97}{source}{}{sourcepointer-5}{}}
\newlabel{sites-cohomology-lemma-internal-hom-evaluate@cref}{{[originallemma][0][]0A97}{[1][0][]source}{}{}{}}
\makeatother
\title{The duality functor of a dualizing complex on an algebraic space}
\author{The Stacks Project Authors}
\date{stacks-project:spaces-duality.tex:local-snapshot-20261012; GNU FDL 1.2 or later}
\begin{document}
\maketitle
\noindent\textit{Editorial note.} All delivered mathematical text is the author's own, copied from the pinned Stacks Project chapter file spaces-duality.tex; no statement or proof step was rewritten, reordered or removed.\par\medskip
\noindent\textit{Editorial note.} The counted proof is the author's complete argument: the reduction to affine etale charts through affine duality, the transfer of coherence and double duality, and the boundedness statements from the ring case.\par\medskip
\noindent\textit{Editorial note.} The label lemma-dualizing-spaces names the counted statement, and the proof environment immediately following it in the source is delivered verbatim as its proof.\par\medskip
\noindent\textit{Editorial note.} The counted statement is carried by the author's own lemma environment in the chapter source.\par\medskip
\noindent\textit{Editorial note.} Cross-chapter same-name check: the label lemma-dualizing-spaces is not carded in any other chapter (the scheme-level labels of duality.tex differ). Pre-open gate: python operations/agent-authoring/tools/precard\_gate.py --source spaces-duality.tex --label lemma-dualizing-spaces --start 390 --end 440 --local-id 28173 returned RESULT=OK.\par\medskip
\noindent\textit{Editorial note.} The statement and proof are copied from the pinned chapter file; the cited inputs are kept as external pointers with their canonical Stacks tags.\par\medskip
\noindent\textit{Allowed standard prerequisites (editorial).}\begin{itemize}
\item Affine duality over an affine etale chart. Same chapter, the lemma comparing the duality functor on the small etale site of an affine scheme with the duality functor on the derived category of the ring. Identifies the duality functor on the space with the algebraic duality functor on each affine etale chart.
\item Coherence and double duality for a dualizing complex over a ring. Dualizing Complexes, the lemma that a dualizing complex over a ring induces an anti-equivalence of the category of coherent complexes with a canonical double duality isomorphism, and its boundedness statements. Provides coherence, the double duality isomorphism, and the exchange of boundedness properties over each affine chart.
\item Restriction of internal hom commutes with localization. Cohomology on Sites, the lemma evaluating the internal hom into a second factor, and the compatibility of internal hom with restriction. Transfers the double duality isomorphism from the charts to the space.
\end{itemize}\medskip
\ifdefined\hypertarget\hypertarget{sourcepointer-1}{}\fi
\noindent\textit{Source reference.} Original lemma 0A86: sources/stacks-chow/tags.txt, source lines 12904--12904. Prerequisite pointer invoked at proof line(s) 413; it names a standard fact of a sibling chapter that the counted argument uses, and every step of the counted proof built on it is retained verbatim in the delivered body.\par\medskip
\ifdefined\hypertarget\hypertarget{sourcepointer-2}{}\fi
\noindent\textit{Source reference.} Original lemma 0A7C: sources/stacks-chow/tags.txt, source lines 12875--12875. Prerequisite pointer invoked at proof line(s) 436; it names a standard fact of a sibling chapter that the counted argument uses, and every step of the counted proof built on it is retained verbatim in the delivered body.\par\medskip
\ifdefined\hypertarget\hypertarget{sourcepointer-3}{}\fi
\noindent\textit{Source reference.} Original reference 0A88: sources/stacks-spaces-duality/pinned-original-source.tex, source lines 307--307. Prerequisite pointer invoked at proof line(s) 414; it names a standard fact of the same chapter that the counted argument uses, and every step of the counted proof built on it is retained verbatim in the delivered body.\par\medskip
\ifdefined\hypertarget\hypertarget{sourcepointer-4}{}\fi
\noindent\textit{Source reference.} Original reference 0A86: sources/stacks-spaces-duality/pinned-original-source.tex, source lines 273--273. Prerequisite pointer invoked at proof line(s) 412; it names a standard fact of the same chapter that the counted argument uses, and every step of the counted proof built on it is retained verbatim in the delivered body.\par\medskip
\ifdefined\hypertarget\hypertarget{sourcepointer-5}{}\fi
\noindent\textit{Source reference.} Original lemma 0A97: sources/stacks-chow/tags.txt, source lines 12940--12940. Prerequisite pointer invoked at proof line(s) 433; it names a standard fact of a sibling chapter that the counted argument uses, and every step of the counted proof built on it is retained verbatim in the delivered body.\par\medskip
\begin{lemma}
\label{lemma-dualizing-spaces}
Let $S$ be a scheme. Let $X$ be a locally Noetherian algebraic space over $S$.
Let $K$ be a dualizing complex on $X$.
Then $K$ is an object of $D_{\textit{Coh}}(\mathcal{O}_X)$
and $D = R\SheafHom_{\mathcal{O}_X}(-, K)$ induces an anti-equivalence
$$
D :
D_{\textit{Coh}}(\mathcal{O}_X)
\longrightarrow
D_{\textit{Coh}}(\mathcal{O}_X)
$$
which comes equipped with a canonical isomorphism
$\text{id} \to D \circ D$. If $X$ is quasi-compact, then
$D$ exchanges $D^+_{\textit{Coh}}(\mathcal{O}_X)$ and
$D^-_{\textit{Coh}}(\mathcal{O}_X)$ and induces an equivalence
$D^b_{\textit{Coh}}(\mathcal{O}_X) \to D^b_{\textit{Coh}}(\mathcal{O}_X)$.
\end{lemma}

\begin{proof}
Let $U \to X$ be an \'etale morphism with $U$ affine. Say $U = \Spec(A)$ and
let $\omega_A^\bullet$ be a dualizing complex for $A$ corresponding to $K|_U$
as in Lemma \ref{lemma-equivalent-definitions} and
Duality for Schemes, Lemma \ref{duality-lemma-equivalent-definitions}.
By Lemma \ref{lemma-affine-duality} the diagram
$$
\xymatrix{
D_{\textit{Coh}}(A) \ar[r] \ar[d]_{R\Hom_A(-, \omega_A^\bullet)} &
D_{\textit{Coh}}(\mathcal{O}_\etale)
\ar[d]^{R\SheafHom_{\mathcal{O}_\etale}(-, K|_U)} \\
D_{\textit{Coh}}(A) \ar[r] &
D(\mathcal{O}_\etale)
}
$$
commutes where $\mathcal{O}_\etale$ is the structure sheaf of the
small \'etale site of $U$. Since formation of $R\SheafHom$ commutes
with restriction, we conclude that $D$ sends
$D_{\textit{Coh}}(\mathcal{O}_X)$ into
$D_{\textit{Coh}}(\mathcal{O}_X)$. Moreover, the canonical map
$$
L \longrightarrow
R\SheafHom_{\mathcal{O}_X}(R\SheafHom_{\mathcal{O}_X}(L, K), K)
$$
(Cohomology on Sites, Lemma \ref{sites-cohomology-lemma-internal-hom-evaluate})
is an isomorphism for all $L$ in $D_{\textit{Coh}}(\mathcal{O}_X)$
because this is true over all $U$ as above by
Dualizing Complexes, Lemma \ref{dualizing-lemma-dualizing}.
The statement on boundedness properties of the functor $D$
in the quasi-compact case also follows from the corresponding
statements of Dualizing Complexes, Lemma \ref{dualizing-lemma-dualizing}.
\end{proof}

\end{document}
